\documentclass{article}
\usepackage{amsmath,amssymb,amsthm,mathtools}
\newtheorem{lemma}{Lemma}
\begin{document}

% Wordings and notation of everyday mathematical writing that are read as the forms of
% docs/writing.md. The last lemma is not read on purpose: `for any` inside a condition
% may ask for one value only.

\begin{lemma}\label{p1}
If $x>0$ then $x^2>0$.
\end{lemma}

\begin{lemma}\label{p2}
Given $x>0$ and $y>0$, we have $xy>0$.
\end{lemma}

\begin{lemma}\label{p3}
For any real number $x$, we have $x^2\ge0$.
\end{lemma}

\begin{lemma}\label{p4}
For any $x\in\mathbb{R}$, $x^2\ge0$.
\end{lemma}

\begin{lemma}\label{p5}
Any real number $x$ satisfies $x^2\geq 0$.
\end{lemma}

\begin{lemma}\label{p6}
For every real number $x$ there is a real number $y$ so that $y>x$.
\end{lemma}

\begin{lemma}\label{p7}
Suppose $x>0$. Then $x^2>0$ holds.
\end{lemma}

\begin{lemma}\label{p8}
Suppose that $x>0$. Then the inequality $x^2>0$ holds.
\end{lemma}

\begin{lemma}\label{p9}
The inequality $x^2+1>0$ holds for all real $x$.
\end{lemma}

\begin{lemma}\label{p10}
Let $x$ be a real number. Then it holds that $x^2\ge0$.
\end{lemma}

\begin{lemma}\label{p11}
Let $x$ be a real number. Then one has $x^2\ge0$.
\end{lemma}

\begin{lemma}\label{p12}
Let $x>0$. Then we obtain $x^2>0$.
\end{lemma}

\begin{lemma}\label{p13}
Let $x>0$. Then there is some $y>0$ with $y<x$.
\end{lemma}

\begin{lemma}\label{p14}
For every $c>0$ there is an integer $b$ with $b\ge c$ such that $b\le c+1$.
\end{lemma}

\begin{lemma}\label{p15}
Let $x$ and $y$ be two positive real numbers. Then $xy>0$.
\end{lemma}

\begin{lemma}\label{p16}
Let $x$ be real. Then $x^2=0$ if and only if $x=0$.
\end{lemma}

\begin{lemma}\label{p17}
Let $x \in \mathbb{R}$. Then $|x| \geqslant 0$ and $|x|=0$ iff $x=0$.
\end{lemma}

\begin{lemma}\label{p18}
Let $f\colon\mathbb{R}\rightarrow\mathbb{R}$ be given by $f(x)=x^2$. Then $f(2)=4$.
\end{lemma}

\begin{lemma}\label{p19}
Let $x>0$. Then $\bigg(\frac{x}{2}\bigg)^2\le x^2$.
\end{lemma}

\begin{lemma}\label{p20}
Let $A\coloneq\{x\in\mathbb{R}:x>0\}$. Then $1\in A$.
\end{lemma}

\begin{lemma}\label{p21}
If $x>0$ for any $x\in\mathbb{R}$, then $1>0$.
\end{lemma}

\end{document}
